\documentclass[11pt]{article}
\usepackage{explorer}
\newcommand{\ExplorerLanguage}{FR}
\newcommand{\ExplorerNumber}{04}
\newcommand{\ExplorerKind}{Solutions détaillées}
\begin{document}
\ExplorerTitle{Mathématiques expertes · approfondissement}{Bézout et équations entières}
\ExplorerSection{Exercice 1}
Calculer PGCD(84,30) et donner une identité de Bézout.\par
\(84 = 2 \times  30 + 24\) et \(30 = 24 + 6\). Donc \(6 = 3 \times  30 - 84\) et le PGCD vaut 6.\par
\ExplorerSection{Exercice 2}
L’équation \(12x + 18y = 5\) admet-elle des solutions entières ?\par
Non : \(PGCD(12,18) = 6\), qui ne divise pas 5.\par
\ExplorerSection{Exercice 3}
Donner toutes les solutions entières de \(5x + 3y = 1\).\par
(−1,2) est une solution. Toutes sont \(x = -1 + 3t, y = 2 - 5t\), avec t entier.\par
\ExplorerSection{Exercice 4}
Parmi les solutions de \(5x + 3y = 31\), trouver celles avec x et y positifs ou nuls.\par
Modulo \(3, x \equiv  2\). Avec \(0 \leq  x \leq  6, x\) vaut 2 ou 5. Les couples sont (2,7) et (5,2).\par
\ExplorerSection{Problème : raisonnement complet}
Une commande de 100 objets doit être livrée en boîtes de 7 et de 11 objets, sans ouvrir de boîte. Déterminer toutes les répartitions possibles, puis celle utilisant le moins de boîtes. Justifier qu’aucune autre solution positive ou nulle n’a été oubliée.\par
On résout \(7x + 11y = 100\). Modulo \(7, 4y \equiv  2\), donc \(y \equiv  4\) (mod 7). Comme \(0 \leq  y \leq  9\), seule la valeur \(y = 4\) convient, d’où \(x = 8\). Il faut 12 boîtes. Cette solution est unique, donc minimise automatiquement le nombre de boîtes.\par
\end{document}
